% NeurIPS 2025 submission template.
% 1) Download neurips_2025.sty from neurips.cc and place it next to this file.
% 2) For anonymous submission: use \usepackage{neurips_2025} (no options).
% 3) For preprint (non-anonymous): \usepackage[preprint]{neurips_2025}.
% 4) For camera-ready (after acceptance): \usepackage[final]{neurips_2025}.

\documentclass{article}

% If you need natbib options, do it before loading neurips_2025:
% \PassOptionsToPackage{numbers,compress}{natbib}

\usepackage{neurips_2025}

\usepackage[T1]{fontenc}
\usepackage[utf8]{inputenc}

\usepackage{microtype}
\usepackage{graphicx}
\usepackage{booktabs}
\usepackage{amsmath,amssymb}
\usepackage{mathtools}
\usepackage{nicefrac}
\usepackage{siunitx}
\usepackage{url}
\usepackage{hyperref}

\usepackage{algorithm}
\usepackage{algpseudocode}

\newcommand{\AMT}{\textsc{amt-ppo}}
\newcommand{\PPO}{\textsc{ppo}}

\title{AMT-PPO: Adaptive Multi-Timescale Trace Memory for Non-Stationary Partially Observable Reinforcement Learning}

% Anonymous submission: keep "Anonymous Author(s)".
\author{Anonymous Author(s)}

\begin{document}

\maketitle

\begin{abstract}
Reinforcement learning agents deployed in the wild must act under partial observability and non-stationarity, process, reward function, or dynamics can drift over time. Recurrent memories (e.g., LSTMs) can address partial observability but often adapt slowly under drift and are difficult to diagnose or control.
We propose AMT-PPO, an on-policy method that augments PPO with an explicit multi-timescale trace memory whose effective forgetting rate is controlled online by a lightweight change statistic computed from streaming error signals (TD residuals and optional feature prediction errors). When change persists, AMT-PPO can trigger internal memory resets and applies a reset-as-truncation variant of generalized advantage estimation (GAE) that preserves valid value bootstrapping while blocking credit assignment across inferred regime boundaries.
We position AMT-PPO as a pragmatic algorithmic component: rather than proving end-to-end convergence under deep function approximation, we provide a principled segmentation semantics, a reference implementation with stabilization for representation drift, and an experimental protocol for partially observable benchmarks with controlled drift.
\end{abstract}


\section{Introduction}
Real-world sequential decision making is rarely stationary and rarely fully observable. Sensors drift, user preferences shift, and reward definitions evolve. Under partial observability, an agent needs memory; under non-stationarity, it also needs mechanisms for fast adaptation and for avoiding negative transfer from outdated information. Standard \PPO with a feedforward policy can fail in partially observable Markov decision processes (POMDPs), while \PPO with recurrent memory can suffer from stale hidden states and slow recovery after distribution shifts.

This paper studies the intersection of (i) online decision making in POMDPs and (ii) unlabeled non-stationarity. We focus on a pragmatic question: can we get most of the benefits of memory and adaptation using a transparent module that is cheaper than recurrence and easier to diagnose?

We propose \AMT, a hybrid of on-policy RL and streaming change detection. The agent maintains a vector memory formed by multiple exponentially decaying traces of learned features. Each trace corresponds to a fixed base timescale (short, medium, long). A drift monitor reads a scalar error signal (temporal-difference residual and optionally a feature prediction error) and produces a normalized change statistic. This statistic controls a gate that increases forgetting when drift is detected, allowing rapid adaptation. When change persists, the method can reset parts of memory to prevent long traces from injecting stale context.

\paragraph{Contributions.}
\begin{itemize}
  \item An explicit multi-timescale trace memory for on-policy RL with drift-conditioned adaptive forgetting.
  \item A simple normalized change statistic based on two exponentially weighted moving averages (EMAs) and an online z-score.
  \item A reset-as-truncation procedure for generalized advantage estimation (GAE) that blocks credit assignment across detected regime boundaries while preserving valid value bootstrapping.
  \item A complete experimental pipeline (controlled drift wrappers, metrics, logging) and a reference implementation.
\end{itemize}

\section{Background}

\subsubsection*{POMDPs}
A POMDP is defined by latent states $s_t \in \mathcal{S}$, actions $a_t \in \mathcal{A}$, observations $o_t \in \mathcal{O}$, transition dynamics $P(s_{t+1}\mid s_t,a_t)$, observation model $P(o_t\mid s_t)$, and rewards $r_t = R(s_t,a_t)$. Since $s_t$ is not observed, a policy typically depends on history or an internal memory state.

\subsubsection*{Non-stationarity}
We consider environments where $P$, $R$, or the observation model can change over time. We do not assume task labels or change-point annotations.

\subsubsection*{PPO and GAE}
\PPO optimizes a clipped surrogate objective for a stochastic policy $\pi_\theta(a\mid m)$. GAE computes advantages using temporal-difference deltas:
\begin{align}
\delta_t = r_t + \gamma V_\phi(m_{t+1}) - V_\phi(m_t),
\end{align}
and aggregates them with a parameter $\lambda$ to trade bias and variance.


\section{Related Work}

Memory for RL in POMDPs is often implemented via recurrence (LSTM/GRU) or attention, and benchmarked on partially observable suites such as POPGym and gridworld navigation. Non-stationary RL is studied under change-point detection, continual RL, and adaptive control. Our method is closest to approaches that introduce explicit forgetting or context tracking, but differs by using a simple, interpretable multi-timescale trace bank and a reset-aware advantage estimator in an on-policy setting.

Our setting combines (i) partial observability, where the agent must act based on histories rather than Markov states, and (ii) non-stationarity, where the environment dynamics and or rewards drift over time. We position AMT-PPO relative to three lines of work: memory for partially observable RL, online adaptation in non-stationary RL, and explicit summary statistics for history.

\paragraph{Reinforcement learning under partial observability.}
Partially observable Markov decision processes (POMDPs) provide the classical formalism for decision-making when observations are incomplete \citep{kaelbling1998planning}. In modern deep RL, the dominant practical approach is to learn a reactive policy augmented with a learnable memory, typically via recurrence. Early deep recurrent agents such as DRQN replace part of a DQN with an LSTM to integrate observations over time \citep{hausknecht2015drqn}. This idea generalizes naturally to actor-critic and policy-gradient methods, where recurrent policies are widely used for partially observable control.

Beyond simple recurrence, distributed RL systems often rely on recurrent policies at scale. For instance, IMPALA uses a decoupled actor-learner architecture and off-policy correction (V-trace), and is frequently paired with recurrent cores for partially observable multitask settings \citep{pmlr-v80-espeholt18a}. R2D2 similarly studies how to train recurrent agents with experience replay, addressing challenges such as representational drift and recurrent-state staleness in distributed settings \citep{DBLP:conf/iclr/KapturowskiOQMD19}. These systems highlight a recurring theme: memory is essential, but its training dynamics can be fragile when representations or data distributions shift.

More recently, attention-based sequence models have become competitive alternatives to RNN memories. GTrXL introduces architectural modifications that stabilize transformers in RL and shows strong performance on memory-demanding tasks \citep{pmlr-v119-parisotto20a}. Transformer-style policies can, in principle, represent richer history dependencies than fixed-size RNN states, but they often come with higher compute and optimization sensitivity. In contrast, AMT-PPO targets a lightweight, explicit memory mechanism that can be dropped into an on-policy PPO-style pipeline without requiring large-scale sequence modeling.

\paragraph{Explicit summaries of history and memory traces.}
A core question in partially observable RL is how to summarize a long observation history into a compact statistic that is both learnable and useful for control. The most common choices (RNN hidden states or attention over contexts) are expressive but often opaque and can be computationally expensive.

\paragraph{Memory traces and explicit history summaries.}
A closely related and very recent direction proposes memory traces as exponential moving averages of history to obtain compact, analyzable representations with theoretical guarantees for value estimation and empirical benefits for control \citep{eberhard2025memorytraces}.
AMT-PPO is philosophically aligned with this approach but targets a different failure mode: in drifting POMDPs, the \emph{appropriate memory timescale itself changes over time}. We therefore (i) maintain a bank of traces at multiple base timescales, and (ii) \emph{control the effective forgetting rates online} using an explicit change statistic, rather than fixing a single decay parameter. This design operationalizes recency-weighting and restart intuitions from non-stationary RL within an on-policy PPO pipeline.


\paragraph{Non-stationary RL: dynamic regret, sliding windows, and adaptation mechanisms.}
A substantial literature studies RL in drifting MDPs using regret or variation-budget frameworks. A representative example is the sliding-window optimism approach, where algorithms such as SWUCRL2-CW and BORL achieve dynamic regret guarantees by restricting estimation to recent data and widening confidence sets to cope with drift \citep{pmlr-v119-cheung20a}. In policy optimization, dynamic regret analyses study how learning tracks a moving target policy and how quickly it can recover when the environment changes \citep{DBLP:conf/nips/FeiYWX20}. These results motivate practical heuristics that implicitly approximate recency weighting and fast resets, but many theoretical methods are designed for tabular or structured settings, while modern deep RL requires function approximation and representation learning.

Recent work has begun to extend non-stationary guarantees beyond tabular models. For instance, \citet{pmlr-v202-feng23e} study non-stationary RL under general function approximation via a dynamic complexity measure. Empirically oriented methods often incorporate explicit change factors or latent context variables. FANS-RL learns factored representations and change factors to improve robustness to non-stationarity \citep{DBLP:conf/nips/FengH0M22}. These approaches are powerful but can require additional modeling assumptions (for example, factorizations or identifiable latent factors) and often rely on auxiliary training objectives.

AMT-PPO instead follows a pragmatic middle route: it does not attempt to explicitly identify latent change causes. It uses a lightweight drift signal to modulate memory timescales and, optionally, perform resets. This connects to a broad class of methods that respond to non-stationarity by emphasizing recent experience (forgetting) or by reinitializing parts of the agent (restarts). Our contribution is to operationalize these ideas inside an on-policy PPO pipeline via an interpretable memory mechanism rather than via a larger policy model.

\paragraph{Meta-RL and fast adaptation.}
Meta-reinforcement learning aims to train agents that adapt quickly to new tasks by learning adaptation strategies across a distribution of tasks. RL$^2$ encodes a fast RL algorithm in an RNN trained by a slower outer RL loop \citep{duan2016rl2}. PEARL performs off-policy meta-RL via probabilistic context variables for task inference and structured exploration \citep{pmlr-v97-rakelly19a}. While meta-RL is conceptually related to non-stationarity, many meta-RL benchmarks assume episodic task boundaries and an underlying task distribution seen during meta-training. In contrast, our focus is on \emph{online} non-stationarity without explicit task segmentation, where the agent must continuously decide how much to remember and how aggressively to update. AMT-PPO can be viewed as a complementary mechanism that provides fast within-lifetime adaptation through explicit, drift-gated memory dynamics, and could potentially be combined with meta-RL objectives.

\paragraph{Summary.}
Across these lines of work, AMT-PPO is most closely related to (i) recurrent and attention-based memory policies for POMDPs \citep{hausknecht2015drqn,pmlr-v119-parisotto20a} and (ii) explicit history summaries such as memory traces \citep{pmlr-v267-eberhard25a}. It is also motivated by non-stationary RL theory and practice emphasizing recency weighting, resets, and adaptation speed \citep{pmlr-v119-cheung20a,DBLP:conf/nips/FeiYWX20}. Our main difference is architectural and algorithmic: we introduce an explicit multi-timescale trace state with drift-adaptive forgetting inside a PPO-style learner, aiming for a simple, interpretable, and compute-efficient mechanism for partially observable, non-stationary online decision-making.


\section{Method: Adaptive Multi-Timescale PPO}

\subsection{Multi-timescale trace memory}
Let $x_t \in \mathbb{R}^d$ be a learned feature embedding of the current observation and previous action. \AMT maintains $M$ traces:
\begin{align}
z_t^{(i)} = (1-\alpha_t^{(i)}) z_{t-1}^{(i)} + \alpha_t^{(i)} x_t,\quad i=1,\dots,M,
\end{align}
where $\alpha_t^{(i)} \in (0,1]$ controls forgetting and is adapted online. We concatenate traces and current policy features to form the policy input memory:
\begin{align}
m_t = \big[\; x_t^{\text{pol}} \;\Vert\; z_t^{(1)} \;\Vert\; \cdots \;\Vert\; z_t^{(M)} \;\big].
\end{align}
Each trace has a base forgetting rate $\alpha_{\text{base}}^{(i)}$ corresponding to an effective window length $L_i \approx 1/\alpha_{\text{base}}^{(i)}$.

\paragraph{Stabilizing the trace feature space.}
A potential failure mode is representation drift: since traces are linear accumulators of feature embeddings, rapidly changing encoders can mix incompatible feature spaces. To mitigate this, we distinguish the memory encoder $f_{\text{mem}}$ used to produce $x_t$ for trace updates from the policy encoder $f_{\text{pol}}$ used to produce $x_t^{\text{pol}}$ (Algorithm~1). In our implementation, $f_{\text{mem}}$ is updated slowly (e.g., as an exponential moving average of $f_{\text{pol}}$ parameters) and we apply normalization to $x_t$ before trace updates. We ablate this design choice in the experiments.


\subsection{Drift signal}
At each step we compute a scalar error $e_t \ge 0$ that correlates with regime change:
\begin{align}
e_t = \lvert \delta_t \rvert + \lambda_{\text{pred}} \, \lVert x_{t+1} - \hat{x}_{t+1} \rVert_2^2,
\end{align}
where $\delta_t$ is a TD residual and $\hat{x}_{t+1}$ is a learned prediction of the next feature. The prediction term is optional and can be disabled ($\lambda_{\text{pred}}=0$).

\subsection{Change statistic and gating}
We maintain two EMAs of the error with different rates:
\begin{align}
\bar{e}_t^{\text{s}} &= (1-\rho_{\text{s}})\bar{e}_{t-1}^{\text{s}} + \rho_{\text{s}} e_t,\\
\bar{e}_t^{\text{l}} &= (1-\rho_{\text{l}})\bar{e}_{t-1}^{\text{l}} + \rho_{\text{l}} e_t,
\end{align}
with $\rho_{\text{s}} \gg \rho_{\text{l}}$. We define a normalized difference:
\begin{align}
d_t = \frac{\bar{e}_t^{\text{s}} - \bar{e}_t^{\text{l}}}{\bar{e}_t^{\text{l}} + \epsilon}.
\end{align}
To make $d_t$ comparable across tasks, we track its EMA mean and variance:
\begin{align}
\mu_t &= (1-\beta)\mu_{t-1} + \beta d_t,\qquad
\sigma_t^2 = (1-\beta)\sigma_{t-1}^2 + \beta (d_t-\mu_t)^2,\\
z_t &= \frac{d_t - \mu_t}{\sqrt{\sigma_t^2 + \epsilon}}.
\end{align}
A soft gate is computed by a sigmoid:
\begin{align}
g_t = \sigma\!\left(\frac{z_t - \tau_{\text{soft}}}{\kappa}\right) \in (0,1),
\end{align}
and interpolates forgetting:
\begin{align}
\alpha_t^{(i)} = \alpha_{\text{base}}^{(i)} + g_t\big(\alpha_{\text{max}}^{(i)} - \alpha_{\text{base}}^{(i)}\big).
\end{align}

\subsection{Hard reset with hysteresis}
A hard reset triggers when $z_t$ exceeds an on threshold $\tau_{\text{on}}$ for $K$ consecutive steps, with hysteresis using an off threshold $\tau_{\text{off}} < \tau_{\text{on}}$ and a cooldown period. Resets can be hard (zero), reset-to-observation, or partial reset of long traces only.

\subsection{Reset-as-truncation for GAE}
We introduce an indicator $\texttt{reset\_done}_t \in \{0,1\}$ marking an internal boundary while the environment continues. We compute deltas using the bootstrap value of the post-reset next memory $m_{t+1}^{\text{post}}$:
\begin{align}
\delta_t = r_t + \gamma \big(1-\texttt{env\_done}_t\big) V_\phi(m_{t+1}^{\text{post}}) - V_\phi(m_t),
\end{align}
and apply truncation in the GAE recursion:
\begin{align}
A_t = \delta_t + \gamma\lambda \big(1-\texttt{env\_done}_t\big)\big(1-\texttt{reset\_done}_t\big) A_{t+1}.
\end{align}

\begin{algorithm}[t]
\caption{\AMT (one environment step)}
\label{alg:amt}
\begin{algorithmic}[1]
\State Input: observation $o_t$, prev-action encoding $a_{t-1}$, traces $\{z_{t-1}^{(i)}\}$, drift statistics
\State Compute features $x_t \leftarrow f_{\text{mem}}(o_t,a_{t-1})$, $x_t^{\text{pol}} \leftarrow f_{\text{pol}}(o_t,a_{t-1})$
\State Compute drift score $z_t$ from $e_t$ via EMAs and z-normalization
\State Gate $g_t \leftarrow \sigma((z_t-\tau_{\text{soft}})/\kappa)$ and interpolate $\alpha_t^{(i)}$
\State Update traces $z_t^{(i)} \leftarrow (1-\alpha_t^{(i)})z_{t-1}^{(i)} + \alpha_t^{(i)} x_t$
\State Form memory $m_t \leftarrow [x_t^{\text{pol}} \Vert z_t^{(1)} \Vert \cdots \Vert z_t^{(M)}]$
\State Sample action $a_t \sim \pi_\theta(\cdot\mid m_t)$, observe $(r_t,o_{t+1})$
\State If hard drift persists: apply reset strategy; set $\texttt{reset\_done}_t \leftarrow 1$
\State Use post-reset $m_{t+1}^{\text{post}}$ for bootstrapped value in GAE
\end{algorithmic}
\end{algorithm}

\noindent\textbf{Remark.} We apply the z-score normalization to the \emph{dimensionless} ratio $d_t$ (Eq.~7), not to the raw error $e_t$, which improves robustness across tasks with different noise

\subsection{Calibration and stability of the drift monitor}
\label{sec:calibration}
The drift monitor introduces a small set of hyperparameters with interpretable time-scale meanings: $(\rho_s,\rho_l)$ control short- vs long-term averaging of the error signal, $\beta$ controls the averaging window for z-normalization, and $(\tau_{\text{soft}},\tau_{\text{on}},\tau_{\text{off}},K)$ control soft gating and persistent hard-reset triggering with hysteresis. To reduce early-training false alarms when TD errors are naturally large, we optionally disable hard resets for an initial warm-up period and rely on hysteresis and cooldown to prevent reset oscillations. We report sensitivity over these parameters and default values in Appendix~A.


\section{Complexity and interpretability}
\AMT adds $O(Md)$ memory state per environment and $O(Md)$ arithmetic per step. Recurrent policies often incur $O(h^2)$ per step for hidden size $h$ and are less transparent. Our internal signals ($z_t$, $g_t$, $\alpha_t^{(i)}$, reset events) enable diagnostics for when and how the agent adapts.

\section{Experiments}

\subsection{Benchmarks}
We recommend:
\begin{itemize}
  \item \textbf{POPGym} tasks that stress memory without heavy visual computation.
  \item \textbf{MiniGrid} partially observable navigation tasks.
\end{itemize}
To study non-stationarity, we apply controlled drift via wrappers: (i) observation shift (additive Gaussian offset with phase-dependent scale), (ii) reward scaling or bias, and optionally (iii) dynamics perturbations where supported.

\subsection{Drift regimes}
\begin{itemize}
  \item \textbf{Piecewise} drift with phases of fixed episode length.
  \item \textbf{Gradual} drift with linear interpolation over a horizon.
\end{itemize}

\subsection{Baselines and ablations}
Baselines: feedforward \PPO with observation stacking, recurrent \PPO (LSTM), \PPO with fixed traces, and traces with reset only. Ablations: drift signal variants (TD only, prediction only, combined), gate parameters and thresholds, reset strategies, and removing truncation in GAE.

\subsection{Metrics}
We report average return, recovery time after change points, reset rate, and compute overhead.

\section{Results}
This section is a placeholder for results produced by the training pipeline. Table~\ref{tab:main} and Figure~\ref{fig:curves} show the intended reporting format.

\begin{table}[t]
\caption{Main results on partially observable non-stationary benchmarks. Fill with mean and standard deviation across seeds.}
\label{tab:main}
\centering
\small
\begin{tabular}{lcccc}
\toprule
Method & Return (mean) & Return (std) & Recovery time & Reset rate \\
\midrule
PPO (stack) & TODO & TODO & TODO & TODO \\
Recurrent PPO (LSTM) & TODO & TODO & TODO & TODO \\
PPO + fixed traces & TODO & TODO & TODO & TODO \\
\AMT (ours) & TODO & TODO & TODO & TODO \\
\bottomrule
\end{tabular}
\end{table}

\begin{figure}[t]
\centering
\fbox{\parbox{0.9\linewidth}{\vspace{2.5cm}\centering Placeholder for learning and recovery curves.\vspace{2.5cm}}}
\caption{Learning curves under drift. Plot median and interquartile range across seeds.}
\label{fig:curves}
\end{figure}

\section{Discussion}
\subsubsection*{When \AMT helps}
\AMT is most relevant when partial observability requires memory and drift is structured enough that the error signal is informative.

\section{Limitations}
AMT-PPO introduces hyperparameters for the drift monitor (EMA rates, thresholds, and persistence). While these parameters have intuitive meanings, they require calibration, and we provide sensitivity analysis and recommended defaults.
Second, the method uses a scalar drift statistic and may miss subtle \emph{localized} shifts that affect only parts of the latent state. Partial resets (e.g., resetting only long traces) mitigate this to some extent, and per-trace or group-wise gating is a natural extension.
Third, trace memory depends on the stability of the feature space used for accumulation. Rapid representation drift in the memory encoder can degrade performance; we mitigate this with a separate slowly-updated memory encoder and normalization, and we ablate these design choices.


\subsubsection*{Failure modes}
If the drift signal is noisy, the gate can oscillate and cause excessive forgetting. If shifts are not reflected in $e_t$, drift may go undetected. If the feature encoder changes too quickly, traces may mix incompatible feature spaces; slow encoder updates or normalization can mitigate this.

\section{Limitations}
AMT-PPO introduces hyperparameters for the drift monitor (EMA rates, thresholds, and persistence). While these parameters have intuitive meanings, they require calibration, and we provide sensitivity analysis and recommended defaults.
Second, the method uses a scalar drift statistic and may miss subtle \emph{localized} shifts that affect only parts of the latent state. Partial resets (e.g., resetting only long traces) mitigate this to some extent, and per-trace or group-wise gating is a natural extension.
Third, trace memory depends on the stability of the feature space used for accumulation. Rapid representation drift in the memory encoder can degrade performance; we mitigate this with a separate slowly-updated memory encoder and normalization, and we ablate these design choices.


\section*{Acknowledgments and disclosure of funding}
\begin{ack}
Do not include acknowledgments in the anonymized submission. In the camera-ready version, include funding sources and competing interests here.
\end{ack}

\bibliographystyle{plainnat}
\bibliography{references}

\appendix

\section{Implementation details}
We use vectorized environments with per-environment trace states, deterministic drift schedules by seeded sampling, separate flags for environment termination and internal reset boundaries, and bootstrapped values computed from post-reset memory for reset-aware GAE.

\section{NeurIPS paper checklist}
NeurIPS requires the checklist to be included after references (and after optional appendices). Fill items with \emph{Yes}, \emph{No}, or \emph{NA}, and include brief justifications.
\begin{enumerate}
  \item \textbf{Claims.} Do you state the main claims and contributions clearly? \textbf{TODO}
  \item \textbf{Limitations.} Do you discuss limitations and failure modes? \textbf{TODO}
  \item \textbf{Reproducibility.} Do you provide sufficient details, code, and hyperparameters? \textbf{TODO}
  \item \textbf{Ethics.} Does the work conform with the NeurIPS code of ethics? \textbf{TODO}
\end{enumerate}

\end{document}
